\clearpage
\fancyhead[L]{\small CERES. PALLAS. JUNO. COMET.}
\SourceLeaf{272}
\begin{center}
CERES, PALLAS, JUNO UND ZWEITER COMET VON 1805.
\end{center}

\noindent\rule{\linewidth}{0.4pt}
\begin{center}
Astronomisches Jahrbuch für 1809. S. 137–140. Berlin 1806.
\end{center}
\noindent\rule{\linewidth}{0.4pt}

\begin{center}
GAUSS AN BODE.
\end{center}
\hfill Braunschweig, den 14. März 1806.

\par\indent Mit Vergnügen theile ich Ew. --- einige im vorigen Winter von mir angestellte Beobachtungen mit, die Ihnen hoffentlich nicht unangenehm sein werden. Sie sind zwar nicht zahlreich; allein das schlechte Wetter, worüber alle meine astronomischen Correspondenten klagen, ist auch hier den ganzen Winter herrschend gewesen: ausserdem kann ich, bei meinen kärglichen Hülfsmitteln, die praktischen Beschäftigungen hier nur als Nebensache betreiben. Zuerst meine Beobachtungen der drei neuen Planeten:

\Needspace{13\baselineskip}
\begin{center}\emph{Ceres.}\end{center}
\begin{center}
\begingroup\fontsize{9}{11}\selectfont\setlength{\tabcolsep}{2pt}\renewcommand{\arraystretch}{1.1}
\begin{tabular}{lr|r|r}
\multicolumn{2}{c|}{Mittlere Zeit} & AR. & Decl.\\ \hline
\textsuperscript{1}805. October 24. & $11^{\mathrm u}\ 22^{\mathrm m}\ 32^{\mathrm s}$ & $109^\circ\ 28'\ 29''$ & $23^\circ\ 27'\ 0''$ Nord.\\
28. & $11\ 31\ 22$ & $110\ 7\ 7$ & $23\ 36\ 58$\\
November 14. & $9\ 22\ 30$ & $111\ 44\ 56$ & $24\ 32\ 1$\\
18. & $9\ 0\ 37$ & $111\ 50\ 48$ & $24\ 48\ 47$\\
\end{tabular}\endgroup
\end{center}

\Needspace{11\baselineskip}
\begin{center}\emph{Pallas.}\end{center}
\begin{center}
\begingroup\fontsize{9}{11}\selectfont\setlength{\tabcolsep}{2pt}\renewcommand{\arraystretch}{1.1}
\begin{tabular}{lr|r|r}
1806. Februar 14. & $8^{\mathrm u}\ 11^{\mathrm m}\ 16^{\mathrm s}$ & $70^\circ\ 16'\ 31''$ & $19^\circ\ 59'\ 13''$ Süd.\\
16. & $7\ 32\ 28$ & $70\ 42\ 39$ & $19\ 20\ 44$\\
17. & $8\ 52\ 38$ & $70\ 56\ 44$ & $19\ 1\ 8$\\
20. & $7\ 49\ 35$ & $71\ 39\ 2$ & $18\ 5\ 0$\\
\end{tabular}\endgroup
\end{center}

\Needspace{8\baselineskip}
\begin{center}\emph{Juno.}\end{center}
\begin{center}
\begingroup\fontsize{9}{11}\selectfont\setlength{\tabcolsep}{2pt}\renewcommand{\arraystretch}{1.1}
\begin{tabular}{lr|r|r}
1806. Februar 17. & $9^{\mathrm u}\ 42^{\mathrm m}\ 0^{\mathrm s}$ & $173^\circ\ 46'\ 45''$ & $0^\circ\ 28'\ 32''$ Nord.\\
20. & $11\ 24\ 15$ & $173\ 15\ 53$ & $0\ 54\ 30$\\
\end{tabular}\endgroup
\end{center}

\par\indent Diesen Beobachtungen zufolge gab meine Ephemeride für die Ceres die ger. Aufst. im Oct. 1805 um 4$\frac{1}{2}'$, im Nov. um 5$\frac{1}{2}'$ zu gross, die Declin. 1$'$ zu klein; für die Pallas die ger. Aufst. um 27$\frac{1}{2}'$ zu gross, die Decl. um 5$'$ zu gross (d. i. zu nördlich); für die Juno die ger. Aufst. um 1$\frac{1}{2}'$ zu gross, die Decl. um 1$'$ zu klein (d. i. zu südlich). [Diese Beobachtungen der Ceres stehen auch Monatl. Corr. B. XIII. S. 189.]



\par\indent Der beträchtliche Unterschied bei der Pallas ist ohne Zweifel Folge der Störungen, die jetzt anfangen sichtbar zu werden: hingegen die über alle Erwartung schöne Übereinstimmung bei der Juno ein Beweis für die Genauigkeit meiner letzten Elemente und meiner Beob. vom 20. Febr. 1805, wonach jene bestimmt waren. Von der Pallas und Juno ist mir bisher gar niemand weiter bekannt, der sie dieses Jahr beobachtet oder gesehen hätte. Pallas ist sehr lichthell und einem Sterne 7., oder 7..8\textsuperscript{ter} Grösse gleich; und muss auch bei voller Erleuchtung sich gut haben im Passage-Instrument und Quadranten beobachten lassen; aber Juno ist freilich schwach an Licht; bei meinen Beobachtungen etwa 11\textsuperscript{ter} Grösse, und daher schwerlich in einem Fernrohre mit Beleuchtung zu beobachten. Ceres ist wie immer gut sichtbar, und ich habe meine Beobachtungen nur deswegen nicht weiter fortsetzen können, weil sie für mein Locale zu hoch kommt. Prof. \textsc{Harding} hat sie öfters gut am Mauerquadranten beobachtet. Mit Vergnügen werde ich Ihnen meine Resultate und den Lauf für die nächste Erscheinung zusenden.

\SourceLeaf{273}

\par\indent Den zweiten Cometen des vorigen Jahres habe ich nur am 8. December zweimal beobachtet, wo er auch dem blossen Auge, etwa mit der Helligkeit eines Sternes 3. oder 4. Grösse gut sichtbar war. Meine Beobachtungen sind diese:
\begin{center}
\begingroup\fontsize{9}{11}\selectfont\setlength{\tabcolsep}{2pt}\renewcommand{\arraystretch}{1.1}
\begin{tabular}{lr|r|r}
\multicolumn{2}{c|}{Mittlere Zeit} & Gerade Aufst. & Decl.\\ \hline
1805. December 8. & $6^{\mathrm u}\ 55^{\mathrm m}\ 17^{\mathrm s}$ & $353^\circ\ 7'\ 40''$ & $23^\circ\ 36'\ 24''$ Südl.\\
& $7\ 58\ 3$ & $352\ 57\ 57$ & $24\ 6$\\
\end{tabular}\endgroup
\end{center}

\par\indent Die parabolischen Elemente, die ich nach dieser Beobachtung, einer \textsc{Olbers}'schen vom 2. Dec. und der im Moniteur bekannt gemachten \textsc{Bouvard}'schen vom 16. Nov. berechnet habe, stehen im Januarheft der M. C. [Seite 266 d. B.]. Allein seitdem habe ich durch Hrn. \textsc{von Zach}'s Güte noch die ganze Reihe der \textsc{Thulis}'schen Originalbeobachtungen erhalten, und eine in Greenwich am 8. Dec. gemachte Meridianbeobachtung ist durch die Zeitungen bekannt geworden. Mit diesen Hülfsmitteln habe ich denn folgende verbesserte Elemente herausgebracht.

\Needspace{11\baselineskip}
\begin{center}
\begingroup\fontsize{8}{10}\selectfont\setlength{\tabcolsep}{2pt}\renewcommand{\arraystretch}{1.1}
\begin{tabular}{lr}
Zeit der Sonnennähe 1805. Dec. 31. . . . . . & $6^{\mathrm u}\ 54^{\mathrm m}\ 35^{\mathrm s}$ Seeberg. M. Z.\\
Länge der Sonnennähe . . . . . . . . . & $109^\circ\ 21'\ 50''5$\\
Länge des aufsteigenden Knotens . . . . . . & $250\ 33\ 34.9$\\
Neigung der Bahn . . . . . . . . . . . & $16\ 30\ 31.9$\\
Logarithmus des kleinsten Abstandes . . . . & $9.9503300$\\
Bewegung . . . . . . . . . . . . . . & rechtläufig.\\
\end{tabular}\endgroup
\end{center}

\par\indent Hier folgt nunmehro die Vergleichung sämmtlicher Beobachtungen mit diesen Elementen. Zwar habe ich noch ein paar andre von \textsc{Thulis} erhalten, allein, da die Sterne womit die Vergleichung gemacht war, in den Sternverzeichnissen und der Hist. Cél. nicht vorkommen, so habe ich diese Beobachtungen nicht reduciren können. [Die Zeit ist für \textsc{Gauss}, \textsc{Olbers}, \textsc{Thulis} bez. die in Braunschweig, Bremen, Marseille.]
\begin{center}
\begingroup\fontsize{8}{10}\selectfont\setlength{\tabcolsep}{2pt}\renewcommand{\arraystretch}{1.1}
\begin{tabular}{r|r|r|r|r|r|l}
1805 & Mittlere Zeit & \shortstack{Berechnete\\gerade\\Aufsteigung} & \shortstack{Berechnete\\Abweichung} & \shortstack{Differ. in der\\ger. Aufst.\\R---B} & \shortstack{Differ. in der\\Abweichung\\R---B} & Beobachter\\ \hline
Nov. 15 & $10^{\mathrm u}\ 28^{\mathrm m}\ 13^{\mathrm s}$ & $14^\circ\ 33'\ 0''$ & $39^\circ\ 18'\ 41''$ N. & $-2'\ 30''$ & $+\phantom{1'}\ 26''$ & Thulis\\
16 & $7\ 17\ 54$ & $14\ 12\ 7$ & $38\ 58\ 53$ & $-1\ 58$ & $+\phantom{1}\ 46$ & Thulis\\
19 & $7\ 11\ 17$ & $12\ 55\ 12$ & $37\ 34\ 27$ & $-\phantom{1}\ 4$ & $+1'\ 7$ & Thulis\\
20 & $7\ 3\ 15$ & $12\ 28\ 23$ & $36\ 59\ 41$ & $-\phantom{1}\ 23$ & $-\phantom{1}\ 32$ & Thulis\\
21 & $9\ 36\ 1$ & $11\ 57\ 5$ & $36\ 16\ 16$ & $+1\ 13$ & $-\phantom{1}\ 20$ & Thulis\\
29 & $6\ 39\ 44$ & $7\ 24\ 39$ & $26\ 47\ 49$ & $+1\ 32$ & $+\phantom{1}\ 1$ & Thulis\\
Dec. 2 & $5\ 34\ 9$ & $4\ 48\ 43$ & $18\ 58\ 18$ & $+1\ 35$ & $-\phantom{1}\ 52$ & Olbers\\
2 & $9\ 39\ 0$ & $4\ 36\ 3$ & $18\ 20\ 28$ & $+1\ 47$ & $-\phantom{1}\ 3$ & Thulis\\
3 & $5\ 48\ 39$ & $3\ 39\ 51$ & $15\ 6\ 38$ & $+\phantom{1}\ 45$ & $-\phantom{1}\ 19$ & Olbers\\
3 & $6\ 30\ 37$ & $3\ 36\ 52$ & $14\ 55\ 11$ & $+1\ 10$ & $-\phantom{1}\ 26$ & Thulis\\
4 & $6\ 31\ 5$ & $2\ 16\ 36$ & $10\ 6\ 29$ N. & $+2\ 12$ & $-1\ 12$ & Thulis\\
6 & $7\ 16\ 16$ & $358\ 39\ 22$ & $3\ 39\ 32$ S. & $+\phantom{1}\ 22$ & $+\phantom{1}\ 48$ & Thulis\\
7 & $7\ 35\ 21$ & $356\ 9\ 21$ & $13\ 4\ 56$ & $\phantom{+1}\ 0$ & $+\phantom{1}\ 42$ & Thulis\\
8 & $5\ 27\ 22$ & $353\ 20\ 30$ & $22\ 56\ 24$ & $-1\ 10$ & $+2\ 44$ & Olbers\\
8 & $6\ 46\ 23$ & $353\ 8\ 33$ & $23\ 34\ 40$ & $-\phantom{1}\ 42$ & $+2\ 18$ & Olbers\\
8 & $6\ 55\ 17$ & $353\ 7\ 44$ & $23\ 35\ 57$ & $+\phantom{1}\ 4$ & $-\phantom{1}\ 47$ & Gauss\\
8 & & $353\ 6\ 21$ & $23\ 40\ 59$ & $-\phantom{1}\ 33$ & $-\phantom{1}\ 15$ & Maskelyne\\
8 & $6\ 59\ 39$ & $353\ 3\ 54$ & $23\ 47\ 34$ & $+\phantom{1}\ 9$ & $-\phantom{1}\ 33$ & Thulis\\
8 & $7\ 58\ 3$ & $352\ 57\ 34$ & $24\ 6\ 39$ & $-\phantom{1}\ 23$ & & Gauss\\
\end{tabular}\endgroup
\end{center}

\par\indent Die Elemente dieses Cometen nähern sich zwar sehr denen des Cometen von 1772, indessen ist doch, besonders bei dem kleinsten Abstande, die Differenz noch zu bedeutend. Eine eigne genaue Un%
\SourceLeaf{274}%
tersuchung der Beobachtungen des Kometen von 1772 hat mir auch gezeigt, dass man die beiderseitigen Elemente einander nicht erheblich näher bringen kann, ohne die Übereinstimmung mit den Beobachtungen zu entstellen. Der Komet von 1772 ist auch keinem Hauptplaneten nahe gekommen, um von dessen Störung eine grosse Änderung seiner Elemente zu leiden. [Diese Bemerkungen theilt \textsc{Gauss} nach \textsc{von Zach} mit unterm 21. Febr. 1806 (Monatl. Corr. B. XIII. S. 310–313. 1806 März.) und gibt zugleich einige vorläufige Andeutungen über Arbeiten, deren Resultate vollständiger in dem weiter unten S. 270 abgedruckten Briefe vom 20. Mai 1806 enthalten sind.]

\par\indent Ich bin jetzt mit meiner Methode, die Planetenbahnen aus den Beobachtungen zu bestimmen beschäftigt und werde in kurzem hoffentlich an die öffentliche Bekanntmachung denken können. Die mannigfaltigen Vervollkommnungen, die ich seit 1801 dabei einzuführen Gelegenheit gehabt habe, machen, dass sie in ihrer Gestalt, nach ihren verschiedenen Theilen ihrer ersten Beschaffenheit fast ganz unähnlich sieht, daher ich es mich nicht gereuen lasse, mit der Publication derselben nicht zu sehr geeilt zu haben.

